\section{Count fluctuations and local one-body covariance}\label{en:consequences}

The retained excitation energy now determines every bounded local
one-body operator compressed by $W_R$ on both sides.
The resulting uniform covariance law is stated in
Proposition~\ref{mom:local-covariance}.  The same energy comparison
controls the cell counts used in the spatial argument.  We record its
four required count and occupation estimates first; the occupation
estimates will follow by applying the covariance law to two projections.
All expectations and dilute errors use the conventions of
Section~\ref{set:section}.

\begin{proposition}[The four energy outputs]\label{en:energy-outputs}
There are arbitrarily small fixed $u>0$ and then $0<w<u/10000$ for which
the following statements hold.  Use the nested partitions fixed in
\eqref{set:nested-lengths}, with $B\asymp D^u$, $R\asymp D^w$, and
$d=1/j$.  For every normalized
nonnegative bosonic ground vector,
\begin{align}
 \frac1{|\mathcal C_R|}\sum_{A\in\mathcal C_R}
       \E(q_A-1)^2&\le C D^{-1-w/6},\label{en:R-count}\\
 \frac1{|\mathcal C_d|}\sum_{A\in\mathcal C_d}
       \E(q_A-1)^4&\le C_jD^{-1-u/6}.\label{en:d-count}
\end{align}
For every normalized, possibly complex, bosonic ground vector,
\begin{align}
 \frac1V\left\langle\sum_{i=1}^N W_{R,i}\right\rangle
       &=\cD\alpha^{3/2}+o(1),\label{en:R-occupation}\\
 \frac1V\left\langle\sum_{i=1}^N W_{d,i}\right\rangle
       &\le \epsilon_j+o_j(1),\qquad \epsilon_j\longrightarrow0
                \quad(j\longrightarrow\infty).\label{en:d-occupation}
\end{align}
Here the thermodynamic limit is taken first, at fixed dilution; the
remainders then tend to zero as $D\to\infty$, and $o_j(1)$ has this meaning
with $j$ fixed.  The bounds are uniform over the indicated ground vectors
and along every density-convergent thermodynamic sequence.  A sufficiently
large volume threshold may depend on the dilution and on $j$.
\end{proposition}

We prove first the corresponding truncated count estimates for arbitrary
ground vectors.  Proposition~\ref{set:count-mgf} will remove the truncation
for nonnegative ground vectors at the end of the proof.  The occupation
estimates themselves do not use positivity.

\subsection{Comparing the upper and lower energies}

We now sum the local estimates.  This step records which positive terms
survive the comparison with \eqref{en:upper}; the following
subsections use them for the count and occupation observables.

On each middle particle-number sector, use the three localization
factors in \eqref{en:three-cutoffs}.  We call the first and third pieces
low pieces, and the second a high piece.  Sums over pieces in what follows
include sector probabilities, spectator integration, and boxes.  We use
the notation $\sum_{\rm low}$ and $\sum_{\rm high}$ for these sums.

\begin{lemma}[The energy surplus]\label{en:surplus}
For $E\asymp Q=D^\theta$, $\theta=u$ or $w$, the following quantities have
sum at most $CVQ^{-1/5}$:
\begin{enumerate}
\item the box-count penalty
\[
 \sum_b Dv_b\,\mathfrak h\left(\frac{n_b}{Dv_b}-1\right)
       \quad\text{in expectation};
\]
\item on middle-sector high pieces, the expectations of
\[
 \sum_{p\ne0}\min(p^2,1)a_p^*a_p
       +D\sum_{A\in\mathcal C_d,\ A\subset b}
          |A|\mathfrak h(q_A-1);
\]
\item on low pieces, the expectations of
$E^{-2}(n_++K_Hn_H)$ and the quantities $\mathcal B(\psi)$;
\item on sectors outside \eqref{en:middle-sector}, their total
$(n_b+Dv_b)$-weighted squared norms.
\end{enumerate}
The first quantity is evaluated on the original vector.  The other
quantities are evaluated on the indicated localized restrictions.
\end{lemma}

\begin{proof}
The preliminary rough bound and \eqref{en:upper} give
\begin{equation}\label{en:preliminary-number}
 \frac1V\sum_b\langle n_+\rangle\le CQ^2D^{1-c_{\mathrm r}}
\end{equation}
on the middle sectors, where $c_{\mathrm r}>0$ is the exponent in
Lemma~\ref{en:rough}.  Lemma~\ref{en:localization} therefore bounds the
sum of localization errors by $CVQ^{-1}$.  In particular this
localization does not consume a fixed fraction of the leading energy.

On a high piece, the support condition $n_L\ge M/8$ gives
\[
 \left\langle\sum_p\min(p^2,1)a_p^*a_p\right\rangle
       \ge cME^{-2}\|\psi\|^2.
\]
Since
$ME^{-2}\asymp DQ^{-399}$, this pays both the rough error
$Cv_bD^{1-c_{\mathrm r}}$ and an $O(v_b)$ LHY scalar, with a fraction tending to
zero, provided $D^{c_{\mathrm r}}Q^{-402}\to\infty$.  The remaining rough bound
then retains a fixed fraction of the kinetic and microcell penalties.
Its leading quadratic scalar, compared with \eqref{en:tangent}, also
retains $cDv_b(n_b/(Dv_b)-1)^2$.  The inequality
$\mathfrak h(s+t)\le C(\mathfrak h(s)+\mathfrak h(t))$ allows a fraction
of this scalar remainder to recenter the microcell penalty from
$n_b/v_b$ to $D$.

On a low piece use Lemma~\ref{en:sharp}.  The convexity of $e_r$ and the
bound $e_r''(x)\ge8\pi r$ imply
\begin{equation}\label{en:scalar-convexity}
 v_be_r(n_b/v_b)-t_b(n_b)
       \ge4\pi\alpha Dv_b\left(\frac{n_b}{Dv_b}-1\right)^2.
\end{equation}
For $n_b<Dv_b/2$, the rough bound itself has a positive difference of
order $Dv_b$ from the tangent; its $D^{1-c_{\mathrm r}}v_b$ error is negligible.
For $n_b>6Dv_b$, divide the labels into
$k=\lfloor n_b/(3Dv_b)\rfloor\ge2$ batches of nearly equal size.
Their densities lie between $3D-o(D)$ and $4.5D+o(D)$; hence they lie
in $[2.5D,5D]$ for large $D$.  Dropping interactions
between batches and applying the rough bound to each batch gives a
linear lower bound whose slope exceeds $e_r'(D)$ by a fixed positive
constant.  Thus it retains $c n_b$ beyond the tangent.  Finitely many
rounding adjustments are harmless for large $Dv_b$.  On these two
classes, $n_b+Dv_b$ also controls every bounded count deviation and the
number of particles to which a bounded one-body observable can apply.

The squared localization factors sum to one, so the scalar tangent sums
to $Ve_r(D)$.  The same identity preserves the box-count penalty, because
box number commutes with all localizations.  Summing the bounds just
obtained and subtracting \eqref{en:upper} proves the assertion, with
$Q^{-1/5}$ in place of the smaller sum of error powers.  The conditions
$u\ll c_*$ and $u\ll c_{\mathrm r}$ ensure that each power error is absorbed.
\end{proof}

\subsection{Second and fourth moments of cell counts}

Apply Lemma~\ref{en:surplus} at scale $E=R$.  On every fixed bounded
interval of relative counts, $\mathfrak h(t)$ is comparable with $t^2$.
Consequently
\begin{equation}\label{en:R-truncated}
 \frac1{|\mathcal C_R|}\sum_{A\in\mathcal C_R}
      \E\min\{(q_A-1)^2,T^2\}
          \le C_TD^{-1}R^{-1/5}.
\end{equation}
Here and below a clipping threshold may be replaced by any fixed
bounded truncation with the same local power behavior.

The fourth moment on a $d$-cell needs more than this scalar convexity.
We return to $E=B$ and use the small number of excitations in a low
piece.  Let $P$ be the projection onto the box constant and let
\[
 X=n_A-n|A|/v_b
      =\sum_{i=1}^n\bigl(\1_A(x_i)-|A|/v_b\bigr).
\]
The one-body operator inside this sum has norm at most one and
vanishing $P$--$P$ entry.

\begin{lemma}[A fourth moment from excitations]\label{en:fourth-block}
On the $n$-particle space, for every vector $\xi$,
\begin{equation}\label{en:fourth-block-bound}
 \|X^2\xi\|^2\le Cn^2\langle\xi,(n_++2)^2\xi\rangle.
\end{equation}
\end{lemma}

\begin{proof}
Decompose the one-body operator into its $P$--$(1-P)$,
$(1-P)$--$P$, and $(1-P)$--$(1-P)$ blocks.  Their second quantizations
change the excitation number by $1,-1,0$, respectively.  Between the
$k$ and $k+1$ sectors the first block has norm at most
$\sqrt{(n-k)(k+1)}$; the diagonal block has norm at most $k$.
Thus $X$ is tridiagonal with block norms bounded by
$C[\sqrt{n(k+1)}+k]$.  The blocks of $X^2$ have bandwidth two and norm
at most $Cn(k+2)$, since $k\le n$.  Squaring, summing the five possible
output blocks, and using Cauchy--Schwarz for their bounded overlaps proves
\eqref{en:fourth-block-bound}.
\end{proof}

On a low piece,
$n_+^2\le C(Mn_++nn_H)$ and $n\le CDQ^3$.
Lemma~\ref{en:surplus} gives
\begin{equation}\label{en:gap-number-bounds}
 \frac1V\sum_{\rm low}\langle n_+\rangle\le CQ^{2-1/5},
 \qquad
 \frac1V\sum_{\rm low}\langle n_H\rangle\le CQ^{2-100-1/5}.
\end{equation}
Divide \eqref{en:fourth-block-bound} by $(D|A|)^4$, average the
$O_j(Q^3)$ microcells in each box, and then average the boxes.  Since
$M\asymp DQ^{-397}$, the result is at most
\begin{equation}\label{en:microcell-powers}
 C_j\left(D^{-1}Q^{-386-1/5}
       +D^{-1}Q^{-86-1/5}+D^{-2}Q^6\right).
\end{equation}
For example, the common prefactor before the summed excitation
moments is $C_jD^{-2}Q^9/V$; this verifies explicitly that averaging
the microcells introduces no missing factor of box volume.
Centering at $D$ rather than $n/v_b$ uses the scalar box-count penalty:
on the middle interval, a bounded fourth power of the relative box
deviation is bounded by a constant times its square.  On high pieces,
the truncated fourth power is bounded by $C\mathfrak h(q_A-1)$, which
is retained in Lemma~\ref{en:surplus}.  The same is true on the
outlying particle-number sectors by their norm penalty.

Finally, for each middle-sector restriction,
$\psi=f^2\psi+h^2\psi$.  Applied to the square root of a nonnegative
truncated count observable, the triangle inequality therefore bounds
its expectation in $\psi$ by twice the sum for these two pieces.
The expectation on the additional $\sqrt2fh$ piece used in the energy
localization is nonnegative and need not be included in this bound.  This establishes the
truncated version of \eqref{en:d-count}; \eqref{en:microcell-powers} is
smaller than $D^{-1}Q^{-1/6}$ for sufficiently small $u$.

For nonnegative ground vectors, Proposition~\ref{set:count-mgf} gives
\begin{equation}\label{en:count-tail-use}
 \mathbb P(n_C>C_0D)\le Ce^{-cD}
\end{equation}
for every unit cell $C$ and a sufficiently large fixed $C_0$.
Hard-core packing bounds all counts in a bounded cell by a fixed power
of $D$.  An $R$-cell contains only polynomially many unit cells; if none
exceeds $C_0D$, its relative count is bounded by $C_0$.  A $d$-cell lies
in one unit cell and has relative count at most $C_0d^{-3}$ on that
event.  Choose these fixed clipping thresholds, depending on $j$ in the
second case.  The omitted second and fourth moments are exponentially
small times a polynomial in $D$.  This removes the truncations and proves
\eqref{en:R-count}--\eqref{en:d-count}, with the slightly weakened
exponents displayed there.

\subsection{Uniform local one-body covariance}

We now extract the one-body information retained by the energy comparison.
On a box $b$ of side $B$, let $W$ be the restriction of $W_R$, and use
local coordinates $[0,B)^3$.  Write $v_b=B^3$ and let $u_p$ be the real
Neumann cosine basis, $p\in(\pi/B)\N_0^3$, defined in
Section~\ref{en:section}.  Define, for $p\ne0$,
\begin{equation}\label{en:limiting-V}
 F_\alpha(p)=\frac12\left\{
       \frac{|p|^2+8\pi\alpha}{\sqrt{|p|^4+16\pi\alpha|p|^2}}-1\right\}.
\end{equation}
The function is integrable on $\R^3$: it is $O(|p|^{-1})$ near zero
and $O(|p|^{-4})$ at infinity, uniformly for $\alpha\in I_\alpha$.

\begin{proposition}[Uniform local covariance]\label{mom:local-covariance}
Use the nested partitions and limiting convention of
Section~\ref{set:section}, with the exponents chosen above.
Let $A=\bigoplus_{b\in\mathcal C_B}A_b$, where the one-particle
operators $A_b$ on $L^2(b)$ are identical in local coordinates and
satisfy $0\le A_b\le W$.  Their cosine matrices may be complex Hermitian.
For every normalized bosonic ground vector,
\begin{equation}\label{mom:local-covariance-formula}
 \frac1V\langle\mathrm d\Gamma(A)\rangle
       =\frac1{v_b}\sum_{p\ne0}(A_b)_{pp}F_\alpha(p)+o(1),
 \qquad (A_b)_{pp}=\langle u_p,A_bu_p\rangle.
\end{equation}
The error is uniform in the ground vector and in all the indicated
operators, which may depend on the scales.  The same formula and
uniformity hold when $A_b=WC_bW$ and $C_b$ is any complex contraction,
identical in local coordinates.  Both assertions hold uniformly for
simultaneous translates of the nested partitions.
\end{proposition}

The proof has three steps.  A matrix estimate converts the quasiparticle
energy of each coherent projection into an error for $A_b$.  The energy
surplus then concentrates a box-volume weighted coherent measure at
density $D$.  Finally this concentration identifies the trace in
\eqref{mom:local-covariance-formula} uniformly in $A_b$.

For $l=R,d$, cellwise Neumann Poincare gives, on a $B$-box,
\begin{equation}\label{en:weight-bounds}
 0\le W_l\le1,\qquad W_l\le Cl^2(-\Delta_{\rm Neu}).
\end{equation}
These operators annihilate the box constant.  In particular every
$0\le A_b\le W$ does so.  No commutation with the Neumann Laplacian
is assumed.

\begin{lemma}[A matrix observable comparison]\label{en:observable}
Let $E=B$, $n$ satisfy \eqref{en:middle-sector}, and $|z|^2\le2n$.
For $0\le A_b\le W$ and an excitation vector $\xi$ in the form domain
of $n_+$, put
\[
 \mathcal Q_z(\xi)=\sum_{p\ne0}\widetilde D_p(z)\|b_p(z)\xi\|^2,
 \qquad
 T_A(z)=\sum_{p\ne0}(A_b)_{pp}|V_p(z)|^2,
\]
where $a_p=U_pb_p+V_pb_p^*$ is the inverse of
\eqref{en:bogoliubov}, with the phase of $z$ restored.
Then $0\le T_A(z)\le Cv_b$ and
\begin{align}
 \left|\langle\xi,\mathrm d\Gamma(A_b)\xi\rangle
                   -T_A(z)\|\xi\|^2\right|
 &\le C(1+R^2)\mathcal Q_z(\xi)\nonumber\\
 &\quad+C\{(1+R^2)\mathcal Q_z(\xi)
                            T_A(z)\|\xi\|^2\}^{1/2}.
 \label{en:observable-bound}
\end{align}
The constant is independent of $A_b$.
\end{lemma}

\begin{proof}
The inverse coefficients satisfy
\begin{equation}\label{en:UV-bounds}
 |U_p|^2=1+|V_p|^2,\qquad
 |V_p|^2=\frac12\left\{
 \frac{\tau(p)+\rho_z\widehat g(p)}{D_p(z)}-1\right\}
       \le C\min(|p|^{-1},|p|^{-4}).
\end{equation}
Rationalizing the expression proves the bound, using
$\rho_z|\widehat g(p)|\le C$, $\tau(p)\asymp p^2$, and positivity
of $\widehat g$ near zero.  On $\mathcal P_L$,
$p^2|U_p|^2\le CD_p(z)$; on its complement $U_p$ is bounded and
$\widetilde D_p(z)\ge cK_H/B^2$.  The same inequalities hold with
$|V_p|$ in place of $|U_p|$.

Write $\mathsf U=\operatorname{diag}(U_p)$,
$\mathsf V=\operatorname{diag}(V_p)$, and
$\mathsf D=\operatorname{diag}(\widetilde D_p)$ on the nonconstant
cosine modes.  Split an input into its low and high frequency parts
before estimating its image under $A_b^{1/2}$.  On the low part use
$A_b\le W\le CR^2\operatorname{diag}(p^2)$; on the high part use
$A_b\le1$ and $K_H/B^2\gg1$.  The squared norm of a sum is at most
twice the sum of squared norms, so the coefficient bounds give
\begin{equation}\label{en:weighted-matrix-bound}
 0\le\mathsf U^*A_b\mathsf U\le C(1+R^2)\mathsf D,
 \qquad
 0\le\mathsf V^*A_b\mathsf V\le C(1+R^2)\mathsf D.
\end{equation}
The split is made before applying $A_b^{1/2}$ because $A_b$ need not
preserve either frequency band.

Take inner products linear in the second variable.  For a finite set
of modes the observable is the squared norm of
$A_b^{1/2}(a_p\xi)_p$.  Decompose this column using
$a_p=U_pb_p+V_pb_p^*$.  The annihilation column has squared norm at
most $q=C(1+R^2)\mathcal Q_z(\xi)$ by
\eqref{en:weighted-matrix-bound}.  For the creation column put
$\mathsf M=\mathsf V^*A_b\mathsf V$.  The canonical commutation
relations give the exact identity
\begin{align}
 \sum_{p,q}\mathsf M_{pq}\langle b_p^*\xi,b_q^*\xi\rangle
 &=\Tr(\mathsf M)\|\xi\|^2
   +\sum_{p,q}(\mathsf M^{\mathsf T})_{pq}
                           \langle b_p\xi,b_q\xi\rangle.
 \label{en:transpose-identity}
\end{align}
The second matrix is the transpose, including when $A_b$ is complex.
Since $\mathsf M$ is Hermitian,
$\mathsf M^{\mathsf T}=\overline{\mathsf M}$, and complex
conjugation preserves its order bounds in
\eqref{en:weighted-matrix-bound}: the dominating matrix $\mathsf D$
is real diagonal.  The creation column therefore has squared norm
$t+r$, where $t=T_A(z)\|\xi\|^2$ and $0\le r\le q$.
Cauchy--Schwarz for the mixed columns gives
\[
 \left|\langle\xi,\mathrm d\Gamma(A_b)\xi\rangle-t\right|
       \le2q+2\sqrt{q(t+q)}\le4q+2\sqrt{qt}.
\]

To bound the trace uniformly, set
$G(p)=\min(|p|^{-1},|p|^{-4})$.  A nonempty lattice shell of upper
radius $t$ contains at most $CB^3t^3$ nonzero cosine frequencies;
nonemptiness implies $Bt\ge\pi$.  Dyadic summation consequently gives
\begin{equation}\label{en:lattice-tails}
 \frac1{v_b}\sum_{0<|p|\le\delta}G(p)\le C\delta^2,
 \qquad
 \frac1{v_b}\sum_{|p|\ge J}G(p)\le C/J
 \quad(0<\delta<1<J),
\end{equation}
and $v_b^{-1}\sum_{p\ne0}G(p)\le C$.  These estimates include
coordinate-plane and axial frequencies.  Since $0\le(A_b)_{pp}\le1$,
we obtain $T_A(z)\le Cv_b$.  The preceding mixed-column estimate
proves \eqref{en:observable-bound}.

The finite-mode computation passes to the full excitation space as
follows.  At fixed parameters $\mathsf U,\mathsf V$ are bounded and
$\sum_p|V_p|^2<\infty$.  The inverse relations and finite
$\langle\xi,n_+\xi\rangle$ imply
$\sum_p\|b_p\xi\|^2<\infty$.  The creation column is square summable
as well, since
\[
 \sum_p|V_p|^2\|b_p^*\xi\|^2
   =\sum_p|V_p|^2\|b_p\xi\|^2
       +\left(\sum_p|V_p|^2\right)\|\xi\|^2<\infty.
\]
Finite coordinate projections commute with $\mathsf U,\mathsf V$
and $\mathsf D$, so the compressed matrices inherit
\eqref{en:weighted-matrix-bound}.  Norm convergence of the columns
and convergence of the scalar traces justify the limit.  If
$\mathcal Q_z(\xi)=\infty$, the asserted bound is automatic.
\end{proof}

\begin{proof}[Proof of Proposition~\ref{mom:local-covariance}]
Apply the sector restrictions and excitation cutoffs used in
Lemma~\ref{en:surplus}, now with $E=B$ and $Q=D^u$.
A box-wise operator evaluates by summing over these restrictions.
For each fixed box the original restriction norms, summed over all
assignments and integrated over spectators, have total square one.
Thus the squared cutoffs preserve both that normalization and the
identity $\sum_b\langle n_b\rangle=N$.

\emph{Reduction to coherent traces.}
On sectors outside \eqref{en:middle-sector}, use
$\mathrm d\Gamma(A_b)\le n$ and the norm penalty in
Lemma~\ref{en:surplus}; their contribution is $o(V)$.
The total one-body localization error on the middle sectors, divided
by $V$, is at most
\[
 CDM^{-2}\le CD^{-1}Q^{794}=o(1).
\]
By splitting the input at $|p|=1$ before applying $W$,
\eqref{en:weight-bounds} also gives
\[
 A_b\le W\le C(1+R^2)\min(-\Delta_{\rm Neu},1).
\]
The high pieces therefore contribute at most
$CV(1+R^2)Q^{-1/5}=o(V)$.

Let $\phi$ be any low piece and $\phi_z$ its zero-mode coherent
projection.  Since $A_b$ annihilates the box constant, coherent
resolution leaves its expectation unchanged.  The part with
$|z|^2>2n$ is negligible by \eqref{en:coherent-tail} and
$\mathrm d\Gamma(A_b)\le n$ on its excitation sectors.  Apply
Lemma~\ref{en:observable} on the remaining region.  The normalization
above gives $\sum_{\rm low}v_b\|\phi\|^2\le V$.  Cauchy--Schwarz
in the piece, spectator, and $z$ variables then bounds the total error
per volume by
\begin{equation}\label{en:occupation-error}
 C(1+R^2)Q^{-1/5}
       +C\{(1+R^2)Q^{-1/5}\}^{1/2}=o(1).
\end{equation}
Here $2w<u/5$, as follows from our choice of exponents.

Define a measure $\nu$ on the low-piece and coherent variables by
\begin{equation}\label{en:coherent-measure}
 \int f\,\dd\nu
  =\frac1V\sum_{\rm low}v_b
       \int_{|z|^2\le2n}f(\phi,z)\|\phi_z\|^2
                                      \frac{\dd^2z}{\pi}.
\end{equation}
As throughout the piece sums, spectator integration is included in
this notation.  The measure is positive and has mass at most one.
Our reduction reads
\begin{equation}\label{en:coherent-trace-reduction}
 \frac1V\langle\mathrm d\Gamma(A)\rangle
       =\int\frac{T_A(z)}{v_b}\,\dd\nu+o(1).
\end{equation}

\emph{Mass and density of the coherent measure.}
The omitted outlying sectors have volume weight at most
$CD^{-1}Q^{-1/5}$.  On a middle-sector high piece,
$n_L\ge M/8$ implies
$\mathrm d\Gamma(\min(p^2,1))\ge cM/B^2$; hence its total omitted
volume weight is bounded by
\begin{equation}\label{en:high-volume-loss}
 \frac1V\sum_{\rm high}v_b\|\phi\|^2
       \le C\frac{v_bB^2}{M}Q^{-1/5}
       \le CD^{-1}Q^{401.8}=o(1).
\end{equation}
The coherent tail removes exponentially small weight.  Therefore
$\nu$ has mass $1-o(1)$.

Here is an explicit form of its density concentration.  Decompose a
fixed $n$-particle piece into excitation-number components $\phi^{(m)}$,
$0\le m\le n$.  Its zero-mode number on that component is $n-m$.
Before restricting to $|z|^2\le2n$, orthogonality of the excitation
sectors shows that the radial coherent law of $s=|z|^2$ is the
mixture of gamma densities
\[
 e^{-s}\frac{s^{n-m}}{(n-m)!}\,\dd s
\]
weighted by $\|\phi^{(m)}\|^2$.  Conditionally on $m$, the mean is
$n-m+1$, the variance is $n-m+1$, and
$\mathbb E|s-(n-m)|\le1+\sqrt{n+1}$.
On the middle range the box penalty controls the squared relative
number deviation, so
\[
 \int\left|\frac n{Dv_b}-1\right|^2\dd\nu
       \le CD^{-1}Q^{-1/5}.
\]
Before the $z$ cut, the box-volume weighted integral of $m/(Dv_b)$ is
\[
 \frac1{DV}\sum_{\rm low}\langle\phi,n_+\phi\rangle
       \le CD^{-1}B^2Q^{-1/5}.
\]
These nonnegative bounds remain valid after the cut.  Combining them
with the conditional gamma estimate yields
\begin{equation}\label{en:coherent-density-concentration}
 \int\left|\frac{\rho_z}{D}-1\right|\dd\nu
 \le C\left(D^{-1/2}Q^{-1/10}
       +D^{-1}B^2Q^{-1/5}+(Dv_b)^{-1/2}\right)=o(1).
\end{equation}
In particular every fixed deviation of $\rho_z/D$ from one has
vanishing $\nu$-weight.

\emph{Identification of the covariance.}
On a fixed annulus $0<\delta\le|p|\le J$, all modes eventually
belong to $\mathcal P_L$, so $\tau(p)-p^2=-\pi/(2B^2)$ there.
The support radius $r$ and \eqref{en:replacement} also imply
$D\widehat g(p)=8\pi\alpha+o(1)$ uniformly on bounded momentum sets.
The expression in \eqref{en:UV-bounds} is continuous in these
parameters on the annulus when $\rho_z/D$ is close to one, with a
denominator bounded away from zero.  By
\eqref{en:coherent-density-concentration},
\[
 \int\sup_{\delta\le|p|\le J}
       \bigl||V_p(z)|^2-F_\alpha(p)\bigr|\,\dd\nu=o(1).
\]
On the exceptional density set use the uniform majorant in
\eqref{en:UV-bounds}.  The same majorant bounds $F_\alpha$, so
\eqref{en:lattice-tails} removes the small and large momentum regions,
uniformly in $A_b$.  Since $|(A_b)_{pp}|\le1$, the annular lattice
count divided by $v_b$ is bounded, and the matrices are identical in
local coordinates, \eqref{en:coherent-trace-reduction} becomes
\eqref{mom:local-covariance-formula}.  All bounds used only
$0\le A_b\le W$ and are therefore uniform in this class.

For a complex contraction $C_b$, write
$C_b=S_b+iT_b$, with $S_b=(C_b+C_b^*)/2$ and
$T_b=(C_b-C_b^*)/(2i)$ Hermitian contractions.  For either $J_b=S_b,T_b$,
\[
 WJ_bW=2W\frac{1+J_b}{2}W-W,
 \qquad 0\le W\frac{1+J_b}{2}W\le W.
\]
A fixed linear combination of the positive-operator identities proves
the assertion for $WC_bW$, preserving uniformity.  Finally, simultaneous
translation of the particles is unitary and preserves the ground space.
Applying the uniform estimate to the translated vector proves the
assertion for every translated nested grid.
\end{proof}

\subsection{The two projection occupations}

It remains to obtain \eqref{en:R-occupation} and
\eqref{en:d-occupation}.  The $d$-cells refine the $R$-cells, so
$P_R\le P_d$ and $0\le W_d\le W_R$.  Both $W_R$ and $W_d$ thus
satisfy Proposition~\ref{mom:local-covariance} on the $B$-boxes.
For $W_R$, averaging a cosine over an $R$-cell gives
$\|P_Ru_p\|\le C/(\delta R)$ uniformly on every annulus
$\delta\le|p|\le J$.  Indeed some component has
$|p_i|\ge\delta/\sqrt3$, and the cosine mean in that coordinate is
bounded by $C/(R|p_i|)$.  Hence $(W_R)_{pp}\to1$ on annuli.
The lattice tails in \eqref{en:lattice-tails} and Neumann Riemann
summation give
\[
 \frac1V\langle\mathrm d\Gamma(W_R)\rangle
   =\frac1{\pi^3}\int_{[0,\infty)^3}F_\alpha(p)\,\dd p+o(1)
   =\int_{\R^3}F_\alpha(p)\frac{\dd p}{(2\pi)^3}+o(1).
\]
The first density is the cosine lattice density; the second equality
uses evenness in each coordinate.  The full integral is
\begin{align}
 \int_{\R^3}F_\alpha(p)\frac{\dd p}{(2\pi)^3}
 &=\frac{(8\pi\alpha)^{3/2}}{4\pi^2}
       \int_0^\infty t^2
       \left\{\frac{t^2+1}{t\sqrt{t^2+2}}-1\right\}\dd t\nonumber\\
 &=\frac{8\alpha^{3/2}}{3\sqrt\pi}.
 \label{en:integral}
\end{align}
The one-dimensional integral equals $\sqrt2/3$, by integrating
$t(t^2+1)/\sqrt{t^2+2}-t^2$ and taking endpoint limits.

For $W_d$, \eqref{en:weight-bounds} gives
$(W_d)_{pp}\le\min(1,Cd^2p^2)$.  The covariance proposition and the
same Riemann sums therefore bound its occupation density by
\[
 \sup_{\alpha\in I_\alpha}\int_{\R^3}
       \min(1,Cd^2p^2)F_\alpha(p)\frac{\dd p}{(2\pi)^3}+o_j(1).
\]
The integral tends to zero as $d=1/j\to0$ by the common integrable
majorant.  Denoting it by $\epsilon_j$ proves
\eqref{en:R-occupation}--\eqref{en:d-occupation} and completes
Proposition~\ref{en:energy-outputs}.
